\chapter{Foundations and Related Work}
\label{ch:foundations}

\section{Document Understanding and Layout Analysis}
\label{sec:document-ai}

\TODO{Concepts before tools: text extraction versus layout recovery versus
semantic structuring. Then the concrete tools this system uses --- PyMuPDF as
the live extractor, with MarkItDown, LlamaParse, MinerU and OpenDataLoader as
swappable text front-ends.}

\section{Large Language Models for Content Generation}
\label{sec:llms}

\TODO{Enough to support \cref{ch:design}: prompting, structured (JSON) output,
batching for cost, and the failure modes that matter here --- hallucination,
prompt injection, and provider unavailability.}

\section{Retrieval-Augmented Generation and Groundedness}
\label{sec:rag}

\TODO{Define groundedness, and distinguish the mechanisms that can enforce it:
scoped retrieval, prompt constraint, refusal gating on a similarity threshold,
and post-hoc citation validation. This taxonomy is what \cref{fig:groundedness}
instantiates, so define the terms here and reuse them.}

\thesisfigure{f3-grounded-vs-ungrounded}
  {Two groundedness regimes over the same corpus. The system contains both,
   with materially different guarantees --- and the stronger one is not the one
   users can reach.}
  {fig:groundedness}

\TODO{\Cref{fig:groundedness} is worth dwelling on because it is a natural
experiment inside a single codebase rather than a comparison across systems:
the same corpus, two tutors, one with a pre-generation refusal gate and one
without. Note that the gate is a pure function, written that way deliberately
so a CI evaluation set can assert on it.}

\section{Spaced Repetition}
\label{sec:srs}

\TODO{SM-2 at the level needed for \cref{ch:implementation}: interval,
ease factor, graduating intervals, and the rating scale. Note that the
implementation hides the algorithm behind a stable interface so FSRS could be
substituted without a schema migration --- a design decision worth citing.}

\section{Gamification in Learning Systems}
\label{sec:gamification}

\TODO{Overview depth only --- this supports RQ2, which is not empirically
answered in this work. Keep it short and resist padding.}

\section{Positioning of this Work}
\label{sec:positioning}

\thesisfigure{f2-positioning}
  {Positioning relative to the surrounding solution space. Document-conversion
   tooling recovers structure without pedagogy; learning platforms provide
   pedagogy but assume authored content; retrieval-augmented question answering
   grounds answers without materialising reusable objects.}
  {fig:positioning}

\TODO{IMPORTANT --- \cref{fig:positioning} is currently derived from this
system's own capabilities and dependency set, NOT from a literature search. It
must not be presented as a survey until you have populated each family with the
specific works you cite in this chapter. Once you do, re-check that the
boundaries still hold; if a cited system spans two families, the framing needs
adjusting rather than the citation being forced.}
